% ============================================================
% Resumen y abstract
% ============================================================
\thispagestyle{empty}
\chapter*{Resumen}
\addcontentsline{toc}{chapter}{Resumen}

El cáncer de pulmón es la primera causa de mortalidad oncológica en el mundo,
con una supervivencia a cinco años que apenas alcanza el 20\,\% y una marcada
heterogeneidad molecular entre subtipos histológicos con implicación
terapéutica directa. Las bases de datos públicas de expresión génica,
principalmente NCBI Gene Expression Omnibus (GEO), reúnen cientos de estudios
transcriptómicos de pulmón, pero su integración manual resulta inviable por
la heterogeneidad de plataformas y por unos metadatos clínicos redactados en
texto libre que dificultan la asignación automática de muestras a grupos
experimentales.

Este Trabajo Fin de Máster presenta \framework, un \emph{framework}
reproducible que automatiza la integración multi-cohorte y la identificación
de biomarcadores transcriptómicos en cáncer de pulmón. El sistema descarga
las cohortes de GEO, las normaliza dentro de cada estudio, delega la
curación de los metadatos clínicos a un modelo de lenguaje de gran tamaño
(Llama~3.3-70b vía Groq), aplica criterios de inclusión, ejecuta análisis
diferencial y entrena clasificadores supervisados (regresión logística
regularizada L1, \emph{Random Forest} y máquinas de vectores soporte) con
validación externa \emph{Leave-One-Dataset-Out} (LODO).

Sobre 11 cohortes GEO descargadas, 8 cumplen los criterios de inclusión
(presencia de casos y controles, curación exitosa y alineamiento
verificado). El framework identifica una firma génica de 1174 genes
replicados en tres cohortes independientes para la tarea de subtipo
histológico (adenocarcinoma frente a carcinoma escamoso), con un panel
mínimo clínicamente manejable de 20 genes que alcanza AUC media LODO de
0,966. La firma completa recupera 18 de 20 marcadores diagnósticos de
inmunohistoquímica clínica sin declararlos: 12 de 12 marcadores escamosos
y 6 de 8 adenocarcinoma. Sobre la tarea tumor \emph{vs.} sano, el modelo
obtiene AUC media LODO 0,925 con balanced accuracy 0,772.

La aportación principal del trabajo es doble: (i) demuestra que la
integración de un modelo de lenguaje en el paso de curación clínica permite
escalar la integración multi-cohorte sin renunciar al rigor metodológico, y
(ii) prueba que la exigencia de replicación en cohortes independientes
—condición ausente en la mayoría de firmas publicadas— produce un panel
biológicamente interpretable y coincidente con marcadores clínicos ya en uso.

\vspace{1cm}
\noindent\textbf{Palabras clave:} transcriptómica, cáncer de pulmón,
adenocarcinoma, carcinoma escamoso, biomarcadores, aprendizaje automático,
modelos de lenguaje, NCBI GEO, validación LODO, firma génica, panel mínimo,
inmunohistoquímica.

\clearpage

\chapter*{Abstract}
\addcontentsline{toc}{chapter}{Abstract}

Lung cancer is the leading cause of cancer-related mortality worldwide, with
a five-year survival rate below 20\,\% and pronounced molecular heterogeneity
between histological subtypes that carries direct therapeutic implications.
Public gene expression repositories, most notably the NCBI Gene Expression
Omnibus (GEO), host hundreds of transcriptomic studies of lung tissue, yet
their manual integration is infeasible due to platform heterogeneity and
free-text clinical metadata that hinder the automatic assignment of samples
to experimental groups.

This Master's Thesis presents \framework, a reproducible framework that
automates multi-cohort integration and transcriptomic biomarker discovery in
lung cancer. The system downloads cohorts from GEO, normalises within each
study, delegates clinical metadata curation to a large language model
(Llama~3.3-70b via Groq), applies inclusion criteria, performs differential
expression analysis and trains supervised classifiers (L1-regularised
logistic regression, Random Forest, Support Vector Machines) with external
Leave-One-Dataset-Out (LODO) validation.

Out of 11 downloaded GEO cohorts, 8 meet the inclusion criteria (presence of
cases and controls, successful curation and verified alignment). The
framework identifies a signature of 1174 genes replicated across three
independent cohorts for the histological subtype task (adenocarcinoma
\emph{vs.} squamous carcinoma), with a minimal clinically manageable panel of
20 genes achieving a mean LODO AUC of 0.966. The full signature recovers
18 of 20 clinical immunohistochemistry markers without being told about them:
12 of 12 squamous markers and 6 of 8 adenocarcinoma markers. On the tumour
\emph{vs.} healthy task, the model attains a mean LODO AUC of 0.925 with
balanced accuracy 0.772.

The main contribution is twofold: (i) it shows that incorporating a large
language model into the clinical curation step allows multi-cohort
integration to scale without sacrificing methodological rigour, and (ii) it
demonstrates that requiring replication across independent cohorts —a
condition absent from most published signatures— yields a biologically
interpretable panel that coincides with clinically established markers.

\vspace{1cm}
\noindent\textbf{Keywords:} transcriptomics, lung cancer, adenocarcinoma,
squamous cell carcinoma, biomarkers, machine learning, large language
models, NCBI GEO, LODO validation, gene signature, minimal panel,
immunohistochemistry.

\clearpage
